\section{A bi-factor scaling law and its identifiability}
\label{sec:bifactor}

\subsection{Model}

Two distinct mechanisms govern the throughput of a permissioned BFT network.
Increasing the number of independent transaction submitters~$\cc$ fills the
transaction pool and the block, raising throughput until a saturation point is
reached. Increasing the validator count~$\n$ adds communication rounds and
signature verification, lowering throughput. We model them multiplicatively.

\begin{Def}[Bi-factor scaling law]
\label{def:bifactor}
For a permissioned BFT network operating under the \RNM{} protocol,
\begin{equation}
\TPS(\n,\cc)=\Tzero\cdot\gfun(\cc)\cdot\n^{-\beta},
\qquad
\gfun(\cc)=\frac{\cc^{\gamma}}{1+\bigl(\cc/\csat\bigr)^{\theta}},
\label{eq:bifactor}
\end{equation}
with $\Tzero>0$, $\beta>0$, $\gamma\in(0,1]$, $\theta>\gamma$, $\csat>0$.
\end{Def}

The constraint $\beta>0$ encodes replicated-state-machine semantics: every
validator executes every transaction, so adding validators cannot increase
aggregate throughput. The constraint $\theta>\gamma$ makes $\gfun$ unimodal,
reproducing the collapse of throughput once the transaction pool overflows. In
the unsaturated regime $\cc\ll\csat$ the model reduces to
$\TPS\approx\Tzero\,\cc^{\gamma}\,\n^{-\beta}$, which is log-linear in both
factors:
\begin{equation}
\log\TPS=\log\Tzero+\gamma\log\cc-\beta\log\n .
\label{eq:loglinear}
\end{equation}

\begin{Remark}[Relation to the Universal Scalability Law]
Expanding~\eqref{eq:usl} for $\kappa\n^{2}\gg1$ gives
$C(\n)\sim(\kappa\n)^{-1}$, \ie{} a power law with exponent $-1$. Model
\eqref{eq:bifactor} therefore contains the coherency-dominated regime of USL as
the special case $\beta=1$, while allowing batching and pipelining to produce
$\beta<1$. We compare both fits empirically in Section~\ref{sec:results}.
\end{Remark}

\subsection{The confounding theorem}

Single-factor studies report $\TPS(\n)\propto\n^{\alpha}$. The following result
explains how such a fit can produce $\alpha>0$ even though $\beta>0$
in~\eqref{eq:bifactor}, and why the reported exponent is not transferable
between testbeds.

The mechanism is simple enough to state without algebra. Suppose an
experimenter doubles the network and, quite reasonably, doubles the load
generator with it. Throughput now moves for two reasons at once: down, because
consensus costs more, and up, because more transactions are being offered. A
regression on $\log\n$ alone has no way to separate the two and attributes the
net movement to $\n$. If the client effect is the stronger of the two, the
fitted exponent comes out positive. Nothing has gone wrong with the protocol;
the regressor was simply not the only thing that changed. What the theorem adds
is the exact bookkeeping: the bias is $\gamma\lambda$, so it is fully determined
by how aggressively the load was scaled and by how much the system benefits from
concurrency.

\begin{Theorem}[Identifiability of the consensus exponent]
\label{thm:identifiability}
Let data $\{(\n_{i},\cc_{i},\TPS_{i})\}_{i=1}^{m}$ be generated
by~\eqref{eq:loglinear} with additive noise $\epsilon_{i}$,
$\E[\epsilon_{i}\mid\n_{i}]=0$, $\Var(\epsilon_{i})=\sigma^{2}<\infty$. Suppose
the measurement campaign follows a \emph{concurrency path}
\begin{equation}
\log\cc_{i}=\lambda\log\n_{i}+\mu+\eta_{i},
\qquad
\E[\eta_{i}\mid\n_{i}]=0,\quad \Var(\eta_{i})<\infty,
\label{eq:path}
\end{equation}
for some $\lambda\in\RR$. Assume further that, writing $x_{i}=\log\n_{i}$, the
empirical moments converge,
\begin{equation}
\bar{x}_{m}\to\bar{x},\qquad
\frac{1}{m}\sum_{i=1}^{m}(x_{i}-\bar{x}_{m})^{2}\ \to\ s_{x}^{2}>0 ,
\label{eq:regularity}
\end{equation}
which holds in particular for a design that revisits a fixed finite set of at
least two validator counts with stable proportions. Let $\ahat_{m}$ be the
ordinary least-squares slope obtained by regressing $\log\TPS_{i}$ on $x_{i}$
alone. Then
\begin{equation}
\ahat_{m}\ \xrightarrow{\ \Prob\ }\ \gamma\lambda-\beta .
\label{eq:confounding}
\end{equation}
In particular $\ahat>0$ if and only if $\lambda>\beta/\gamma$, and $\ahat=-\beta$
if and only if $\lambda=0$, \ie{} only when the client concurrency is held fixed.
\end{Theorem}

\begin{Proof}
Write $x_{i}=\log\n_{i}$ and $y_{i}=\log\TPS_{i}$. Substituting~\eqref{eq:path}
into~\eqref{eq:loglinear},
\begin{equation}
y_{i}
=\log\Tzero+\gamma\bigl(\lambda x_{i}+\mu+\eta_{i}\bigr)-\beta x_{i}+\epsilon_{i}
=\underbrace{\bigl(\log\Tzero+\gamma\mu\bigr)}_{\eqdef\,a}
+\bigl(\gamma\lambda-\beta\bigr)x_{i}
+\underbrace{\gamma\eta_{i}+\epsilon_{i}}_{\eqdef\,\xi_{i}} .
\label{eq:reduced}
\end{equation}
Equation~\eqref{eq:reduced} is a correctly specified simple linear regression
with slope $\gamma\lambda-\beta$ and disturbance $\xi_{i}$. By hypothesis
$\E[\xi_{i}\mid x_{i}]=\gamma\E[\eta_{i}\mid\n_{i}]+\E[\epsilon_{i}\mid\n_{i}]=0$,
and $\Var(\xi_{i})\le2\gamma^{2}\Var(\eta_{i})+2\sigma^{2}<\infty$.

The least-squares slope is the ratio of empirical moments
\begin{equation*}
\ahat_{m}
=\frac{\frac{1}{m}\sum_{i}(x_{i}-\bar{x}_{m})(y_{i}-\bar{y}_{m})}
      {\frac{1}{m}\sum_{i}(x_{i}-\bar{x}_{m})^{2}}
=\bigl(\gamma\lambda-\beta\bigr)
+\frac{\frac{1}{m}\sum_{i}(x_{i}-\bar{x}_{m})(\xi_{i}-\bar{\xi}_{m})}
      {\frac{1}{m}\sum_{i}(x_{i}-\bar{x}_{m})^{2}} ,
\end{equation*}
where the second equality follows by substituting~\eqref{eq:reduced} and noting
that the constant $a$ cancels against the centring. For the numerator, the
summands $(x_{i}-\bar{x})\xi_{i}$ have zero mean by the law of iterated
expectations, $\E[(x_{i}-\bar{x})\xi_{i}]=\E[(x_{i}-\bar{x})\E[\xi_{i}\mid x_{i}]]=0$,
and finite variance; the weak law of large numbers therefore gives
\begin{equation*}
\frac{1}{m}\sum_{i=1}^{m}(x_{i}-\bar{x}_{m})(\xi_{i}-\bar{\xi}_{m})
\ \xrightarrow{\ \Prob\ }\ 0 .
\end{equation*}
The denominator converges to $s_{x}^{2}>0$ by~\eqref{eq:regularity}. Applying
the continuous mapping theorem to the ratio yields~\eqref{eq:confounding}.

For the equivalences, $\gamma>0$ by Definition~\ref{def:bifactor}, so
$\gamma\lambda-\beta>0\iff\lambda>\beta/\gamma$, and
$\gamma\lambda-\beta=-\beta\iff\gamma\lambda=0\iff\lambda=0$.
\end{Proof}

\begin{Remark}[Order of magnitude]
The bias is not a second-order correction. With $\gamma=0.8$, $\beta=0.5$ and a
load generator scaled proportionally to the network ($\lambda=1$), the theorem
gives $\ahat=+0.3$: the same measurements that contain a genuine $-0.5$ decay
report growth. Extrapolated from $\n=5$ to $\n=60$ the two exponents differ by
a factor $12^{0.8}\approx7.3$, which is the order of the discrepancy we set out
to explain.
\end{Remark}

\begin{Corollary}[Non-transferability]
\label{cor:nontransfer}
The single-factor exponent $\ahat$ is a property of the pair (protocol,
measurement design). Two laboratories measuring the same implementation with
different concurrency paths $\lambda_{1}\ne\lambda_{2}$ obtain exponents
differing by $\gamma(\lambda_{1}-\lambda_{2})$, and predictions extrapolated to
$\n_{t}$ differ by the factor $\n_{t}^{\gamma(\lambda_{1}-\lambda_{2})}$.
\end{Corollary}

\begin{Corollary}[Sign reversal]
\label{cor:sign}
If a load generator is scaled proportionally to the network, $\cc\propto\n$
so $\lambda=1$, then $\ahat=\gamma-\beta$, which is positive whenever
$\gamma>\beta$. A single-factor fit then reports throughput \emph{growing} with
the validator count, contradicting replicated-state-machine semantics.
\end{Corollary}

Corollary~\ref{cor:sign} is the mechanism we conjecture behind previously
reported positive exponents, including our own~\cite{peniak2026a}.
Section~\ref{sec:results} tests it directly by estimating $\gamma$, $\beta$ and
$\lambda$ separately and checking the identity~\eqref{eq:confounding}.

\subsection{Estimation}

Given measurements $\{(\n_{i},\cc_{i},\TPS_{i})\}_{i=1}^{m}$ acquired under the
\RNM{} protocol on a factorial grid, the unsaturated part
($\cc_{i}\le\csat$) is fitted by ordinary least squares in log-space,
\begin{equation}
\min_{\log\Tzero,\gamma,\beta}\ \sum_{i=1}^{m}
\Bigl(\log\TPS_{i}-\log\Tzero-\gamma\log\cc_{i}+\beta\log\n_{i}\Bigr)^{2},
\label{eq:ols}
\end{equation}
which is a two-regressor linear problem solvable in closed form. The saturation
parameters $(\csat,\theta)$ are then obtained by nonlinear least squares
on the full $\cc$-sweep with $(\Tzero,\gamma,\beta)$ held fixed. A factorial
design guarantees $\Cov(\log\n,\log\cc)=0$, so the two effects are orthogonal by
construction and Theorem~\ref{thm:identifiability} applies with $\lambda=0$.

\begin{Algorithm}[!t]
\caption{Bi-factor calibration under the \RNM{} protocol}
\label{alg:calibration}
\REQUIRE Grid $\mathcal{N}\times\mathcal{C}$, quota $\q$, replicate count $r$
\ENSURE $\widehat{\Tzero},\ghat,\bhat$, saturation $(\widehat{\csat},
        \widehat{\theta})$, residual scale $\hat\sigma$
\STATE \textbf{for} each $(\n,\cc)\in\mathcal{N}\times\mathcal{C}$ and each of
       $r$ replicates: deploy $\n$ validators with quota $\q$, drive $\cc$
       independent senders, record steady-state $\TPS$
\STATE Discard warm-up windows; aggregate replicates by median
\STATE Select the unsaturated subset $\mathcal{U}=\{i:\TPS_{i}\ \text{increasing
       in}\ \cc\}$
\STATE Solve~\eqref{eq:ols} on $\mathcal{U}$ $\Rightarrow$
       $\widehat{\Tzero},\ghat,\bhat$
\STATE Fit $(\widehat{\csat},\widehat{\theta})$ by nonlinear least squares on
       the full $\cc$-sweep
\STATE $\hat\sigma\leftarrow$ standard deviation of log-space residuals
\STATE \RET $\widehat{\Tzero},\ghat,\bhat,\widehat{\csat},\widehat{\theta},\hat\sigma$
\end{Algorithm}
